\documentclass[11pt]{article}

\usepackage[a4paper,margin=1in]{geometry}
\usepackage{fancyhdr}
\usepackage{lastpage}
\usepackage{titlesec}
\usepackage{booktabs} 
\usepackage{enumitem} 
\usepackage{hyperref}
\usepackage{graphicx}
\usepackage{amsmath}
\usepackage{xcolor}
\usepackage{listings}

% Dark Theme Colors
\definecolor{pageblack}{rgb}{0,0,0}
\definecolor{textgrey}{rgb}{0.5,0.5,0.5}
\definecolor{codegreen}{rgb}{0,0.6,0}
\definecolor{codegray}{rgb}{0.5,0.5,0.5}
\definecolor{codepurple}{rgb}{0.58,0,0.82}
\definecolor{backcolour}{rgb}{0.05,0.05,0.05} % Very dark grey for code blocks

% Apply Dark Theme
\pagecolor{pageblack}
\color{textgrey}

% Setup listings for CUDA/C++
\lstset{
    language=C++,
    backgroundcolor=\color{backcolour},   
    commentstyle=\color{codegreen},
    keywordstyle=\color{magenta},
    numberstyle=\tiny\color{codegray},
    stringstyle=\color{codepurple},
    basicstyle=\ttfamily\footnotesize\color{textgrey},
    breakatwhitespace=false,         
    breaklines=true,                 
    captionpos=b,                    
    keepspaces=true,                 
    numbers=left,                    
    numbersep=5pt,                  
    showspaces=false,                
    showstringspaces=false,
    showtabs=false,                  
    tabsize=2
}

% Header and Footer Setup
\pagestyle{fancy}
\fancyhf{}
\lhead{\color{textgrey}Project 4}
\chead{\color{textgrey}CUDA Optimization}
\rhead{\color{textgrey}Semester: SPRING 2026}
\lfoot{\color{textgrey}Arika Khor}
\cfoot{}
\rfoot{\color{textgrey}Page \thepage\ of \pageref{LastPage}}

\begin{document}

\begin{center}
    {\Large \textbf{\color{white}Project 4 Report: Tree Reduction and Stencil Computation in GPU}}\\
    \vspace{0.5em}
    Author: Arika Khor
\end{center}

\vspace{1em}

\section*{\color{white}Problem 1 \& 2: PDNcoin Mining and Reduction}

\subsection*{Algorithm Description}
The goal of these problems is to optimize a cryptocurrency mining workflow using CUDA. The process is divided into three parallelizable stages:
\begin{enumerate}[label=\arabic*.]
    \item \textbf{Nonce Generation:} Embarrassingly parallel generation of random trial values.
    \item \textbf{Hash Calculation:} Compute-intensive modular arithmetic per nonce.
    \item \textbf{Global Minimum Search:} Finding the winning nonce with the lowest hash value.
\end{enumerate}

In \textbf{Problem 2}, we implemented a high-performance tree reduction to replace the CPU bottleneck. The primary technical challenge was returning both the minimum hash and its associated nonce. We solved this by packing both 32-bit values into a single 64-bit \texttt{unsigned long long} and using \texttt{atomicCAS} for global synchronization.

\subsection*{Optimized Reduction Kernel Snippet}
\begin{lstlisting}
__global__ void reduction_kernel(unsigned int *d_hashes, unsigned int *d_nonces, int num_elements, unsigned long long *d_min_result) {
    extern __shared__ unsigned int s_data[];
    unsigned int *s_hashes = s_data;
    unsigned int *s_nonces = &s_data[blockDim.x];

    int tid = threadIdx.x;
    int i = blockIdx.x * blockDim.x + threadIdx.x;

    if(i < num_elements) {
        s_hashes[tid] = d_hashes[i];
        s_nonces[tid] = d_nonces[i];
    } else {
        s_hashes[tid] = 123123123; // MAX
        s_nonces[tid] = 0xFFFFFFFFU;
    }
    __syncthreads();

    for (unsigned int s=blockDim.x/2; s>0; s>>=1) {
        if (tid < s) {
            if (s_hashes[tid] > s_hashes[tid+s]) {
                s_hashes[tid] = s_hashes[tid+s];
                s_nonces[tid] = s_nonces[tid+s];
            }
        }
        __syncthreads();
    }

    if (tid == 0) {
        unsigned long long local_min = ((unsigned long long)s_hashes[0] << 32) | (unsigned long long)s_nonces[0];
        atomicMin(d_min_result, local_min); // Logic handled via atomicCAS in implementation
    }
}
\end{lstlisting}

\subsection*{Performance Results (Schooner Cluster)}

\subsubsection*{Runtime (seconds) - Mining Trials}
\begin{center}
\begin{tabular}{@{}lcc@{}}
\toprule
Implementation & 5 Million Trials & 10 Million Trials \\
\midrule
\textbf{Problem 1 (CPU Reduction)} & 0.678518 & 1.131965 \\
\textbf{Problem 2 (GPU Reduction)} & 0.632880 & 1.087014 \\
\bottomrule
\end{tabular}
\end{center}

\noindent\rule{\textwidth}{0.4pt}

\section*{\color{white}Problem 3 \& 4: Tiled Convolution and Max-Pooling}

\subsection*{Algorithm Description}
These problems implement a deep learning inference pipeline (Stencil Computation). We implemented a 2D convolution layer followed by a max-pooling operation.

Key optimizations include:
\begin{enumerate}[label=\arabic*.]
    \item \textbf{Shared Memory Tiling:} Each thread block collaboratively loads an input tile plus a 2-pixel halo into shared memory, reducing global memory pressure.
    \item \textbf{Dynamic Configuration:} The kernels use \texttt{extern __shared__} to support runtime-configurable \texttt{BLOCK\_SIZE} and \texttt{RADIUS} arguments.
    \item \textbf{Kernel Chaining:} Problem 4 pipes the convolution output directly to the max-pooling kernel while staying entirely on the device.
\end{enumerate}

\subsection*{Performance Results (2048 x 2048 Matrix)}
\begin{center}
\begin{tabular}{@{}lc@{}}
\toprule
Pipeline Stage & Runtime (seconds) \\
\midrule
\textbf{Convolution Only} & 0.001332 \\
\textbf{Convolution + Max-Pooling} & 0.002001 \\
\bottomrule
\end{tabular}
\end{center}

\noindent\rule{\textwidth}{0.4pt}

\section*{\color{white}AI Usage and Workflow}

\subsection*{Gemini CLI \& GSD Methodology}
This project utilized the **Gemini CLI** integrated with the **GSD (getshitdone)** workflow for rigorous planning and execution.

\begin{itemize}
    \item \textbf{Initialization:} Automatically mapped the project structure from the PDF instructions and established coding standards (e.g., \texttt{gpuErrchk} requirement).
    \item \textbf{Research:} Analyzed the provided serial code to identify optimization targets, specifically pinpointing the 40MB D2H transfer in Problem 1 as a primary bottleneck.
    \item \textbf{Execution:} Used AI-driven surgical edits to migrate static shared memory allocations to dynamic pointers, enabling the configurable CLI requirement without manual boilerplate.
    \item \textbf{Verification:} Generated a Google Test suite to stress-test kernels with non-multiple grid sizes and edge-case radii.
\end{itemize}

\noindent\rule{\textwidth}{0.4pt}

\section*{\color{white}Project Layout}
\begin{itemize}
    \item \textbf{common/}: Centralized timers and error handling (\texttt{gpuErrchk}).
    \item \textbf{Problem\_1/}: GPU mining with host-side reduction.
    \item \textbf{Problem\_2/}: Fully optimized 64-bit atomic GPU reduction.
    \item \textbf{Problem\_3/}: Tiled convolution with dynamic halo loading.
    \item \textbf{Problem\_4/}: Chained GPU pipeline (Conv + MaxPool).
    \item \textbf{tests/}: GTest unit testing infrastructure.
\end{itemize}

\end{document}
